\subsection{Cohorts}

The study is built on seven public 12-lead cohorts spanning five countries plus
one restricted institutional cohort, selected so that the transfer matrix
contains the contrasts the theory needs rather than merely many rows.

\begin{table}[t]
\centering\small
\begin{tabular}{llrll}
\toprule
Cohort & Country & Records & Labels from & Role \\
\midrule
MIMIC-IV-ECG      & US & 800{,}035 & cart statements      & source domain \\
Georgia 12-lead   & US & 10{,}344  & SNOMED-CT            & \textbf{same-country control} \\
PTB-XL            & DE & 21{,}799  & cardiologist SCP     & target \\
Chapman-Shaoxing  & CN & 10{,}247  & SNOMED-CT            & target \\
Ningbo            & CN & 34{,}905  & SNOMED-CT            & target \\
CPSC-2018         & CN & 6{,}877   & SNOMED-CT            & target \\
CODE-15\%         & BR & 345{,}779 & validated automatic  & target \\
EchoNext          & US & $\sim$100{,}000 & \emph{echocardiography} & label-provenance control \\
Kangbuk Samsung   & KR & ---       & over-read + screening & \textbf{prevalence contrast} \\
\bottomrule
\end{tabular}
\caption{Cohorts. Two are controls rather than additional targets. Georgia
shares a country with the source but not a health system, so it separates a
national effect from a generic site effect. EchoNext takes its labels from a
separate imaging modality, so a cliff that persists there cannot be an artefact
of how sites read tracings.}
\label{tab:cohorts}
\end{table}

\paragraph{Why a Korean screening cohort.}
Every publicly available ECG cohort is a hospital population, and
hospital-to-hospital transfer confounds the two shifts the decomposition
separates: sites that see more disease also see different disease.  An
asymptomatic health-screening population breaks the confounding --- prevalence
falls by an order of magnitude while the signal-generating process is
unchanged.  It is therefore the cohort on which $D_{\mathrm{label}}$ and
$D_{\mathrm{concept}}$ are most nearly orthogonal, and the one that makes
Theorem~\ref{thm:decomp} empirically identifiable rather than merely true.

\paragraph{Federated handling.}
The Korean cohort cannot leave its institution, and it does not need to.  Every
estimator in this paper depends on the target site only through per-bin counts,
mean scores and event rates, plus nonconformity order statistics.  The model
and the source bin edges travel inbound; roughly 1.5\,kB of aggregates travels
out.  A disclosure audit refuses any payload containing a cell describing fewer
than ten people.

\subsection{Label harmonisation}

CODE-15 is annotated for six rhythm and conduction abnormalities, which bounds
any label set shared by every cohort.  We therefore define \textsc{core6}
(atrial fibrillation/flutter, first-degree AV block, complete right and left
bundle branch block, sinus bradycardia, sinus tachycardia) for the headline
transfer matrix, and \textsc{ext12}, adding six morphology and hypertrophy
labels, for the PhysioNet-family cohorts.

One rule in the harmonisation does disproportionate work: a label a cohort does
not annotate is recorded as \emph{unobserved}, never as negative.  Treating an
annotation gap as a confirmed absence manufactures a prevalence difference out
of nothing, and that artefact would land squarely in the label-shift component
this study exists to measure.  Every label matrix therefore carries an
observation mask, and masked entries contribute to neither training nor
evaluation.

\subsection{Preprocessing}

Identical for every cohort and applied in a fixed order: strip the zero padding
CODE-15 ships, band-pass 0.5--47\,Hz with a zero-phase filter, resample to
250\,Hz, centre-crop to 10\,s, and scale each lead by its median and
interquartile range.  Filtering precedes resampling so the anti-alias behaviour
is set by the declared low-pass rather than by the resampler, making the result
independent of a cohort's native rate.

A mains notch is deliberately \emph{not} applied.  Mains is 50\,Hz in Germany
and China and 60\,Hz in the Americas and Korea, so a per-cohort notch would
remove a genuine site signature, improving apparent transfer while changing
nothing a deployed model would face.

\subsection{Model and adaptation protocol}

The encoder is a 1-D residual network over the twelve leads with global average
pooling, sized so that every experiment runs on CPU.  The protocol is the
foundation-model one:

\begin{enumerate}[leftmargin=*,itemsep=1pt]
\item \textbf{Pretrain} by masked-patch reconstruction on the pooled
\emph{unlabeled} union of all sites.  Contrastive pretraining was rejected for
a reason specific to this study: the standard ECG augmentation set makes the
encoder invariant to exactly the acquisition differences that distinguish these
sites, which would erase the effect by construction.
\item \textbf{Probe} on source labels only, with the encoder frozen.  The head
is strictly linear, so a transfer failure cannot be attributed to head
capacity.  It is fitted to convergence rather than by a few gradient steps ---
with calibration as the endpoint, an undertrained head is miscalibrated for
reasons unrelated to shift.
\item \textbf{Calibrate} by temperature scaling on a held-out source split,
establishing Assumption~\ref{ass:srccal}.
\item \textbf{Freeze} and score every site.  No target labels, no target
statistics and no per-site renormalisation touch the model again; the frozen
backbone is held in evaluation mode so that even batch-normalisation running
statistics cannot adapt.
\end{enumerate}

Splits are at the patient level throughout, with a leakage assertion executed
rather than assumed.  Record-level splitting inflates discrimination, and this
paper's headline is that discrimination is \emph{preserved}: leakage would
manufacture its own result.

\subsection{Simulation}

The theorems are statements about estimators --- that a decomposition is
identified, that a label count suffices, that a coverage guarantee survives a
bounded perturbation --- and claims of that kind cannot be checked on real
cohorts, where the ground-truth split between label and concept shift is
precisely what is unknown.  We therefore also run the entire pipeline on a
generator with known $(\pi_t,\eta)$.

Beats are synthesised as a three-dimensional cardiac dipole and projected onto
the twelve leads through an inverse-Dower matrix, so a conduction abnormality is
introduced once and appears with physiologically consistent polarity across all
leads; the derived limb leads are computed from I and II as a real recording
cart does, so Einthoven's and Goldberger's relations hold exactly.  Prevalence,
morphology and acquisition are controlled independently, matching the three
things that actually differ between these cohorts.  Prevalences are set near
published figures for the corresponding real cohorts, including an order-of-
magnitude lower screening arm.

\subsection{Estimation}

Calibration is reported as $L^1$ ECE with 15 equal-mass bins and as the
debiased $L^2$ estimator, which subtracts the per-bin binomial variance.  The
debiased value is not clipped at zero: clipping reintroduces positive bias
exactly in the regime where one would claim a model is well calibrated.
Equal-mass rather than equal-width bins are used because the score distribution
piles up near zero at a screening site, which is where equal-width binning
fails.  Intervals are stratified bootstraps, stratified on the label so that
replicates cannot lose all the positives.  Any per-label estimate resting on
fewer than 25 positive cases is reported but flagged.
